\vspace*{-1.2cm}

\begin{center}
  {\LARGE \textbf{Simone Romiti}}\\[2pt]
  {\large Curriculum Vitae}
\end{center}

\vspace{2pt}

\cvhead{Personal information}

\noindent
\begin{tabularx}{\textwidth}{@{\hspace{0.5cm}}p{8.2cm}X@{}}
\textit{Simone Romiti} & Institute for Theoretical Physics \\
Date of birth: 25/08/1994 & Albert Einstein Center for Fundamental Physics \\
Place of birth: Rome, Italy & University of Bern \\
Nationality: Italian & Sidlerstrasse 5 \\
Email: \href{mailto:simone.romiti.1994@gmail.com}{simone.romiti.1994@gmail.com} & 3012 Bern, Switzerland \\
Phone: +39 334 356 3288 & \\
Website: \href{https://simone-romiti.github.io/}{simone-romiti.github.io} & \\
\end{tabularx}

\cvhead{Research activity}

My field of research is lattice field theory, a first-principles approach to the non-perturbative regime of quantum field theories.
I work in both of its formulations: the Lagrangian one, where observables are computed by Monte Carlo sampling of Euclidean path integrals, and the Hamiltonian one, which is used, among others, by tensor-network and quantum-computing methods.
I have made contributions to:

\begin{tightitemize}
  \item the solution of the eigenvalue problem of $\mathrm{SU}(N)$ lattice Hamiltonians with physics-informed
  neural networks. In a sole-author paper I introduced an adiabatic training scheme that follows eigenfunctions
  and eigenvalues along the coupling dependence without an explicit truncation of the Hilbert space
  \ref{pub:pinn}. This is the method that the present proposal develops further.
  \item the digitisation of gauge fields for quantum simulation. I worked on discrete formulations of the canonical momenta in $\mathrm{SU}(2)$ using subsets of the group manifold \ref{pub:canonicalmomenta}.
  I introduced a discrete transform between the points on the sphere $S_3$ and their momenta, minimising the breaking of the canonical commutation relations \ref{pub:digitization}.
  \item the quantitative matching of the Lagrangian and the Hamiltonian formulations at finite lattice spacing,
  where they differ by discretisation effects \ref{pub:matching}, and its application to the determination of the running coupling
  of $(2+1)$-dimensional quantum electrodynamics with Monte Carlo and quantum-computing methods \ref{pub:runningcoupling}.
  \item the lattice input to the anomalous magnetic moment of the muon within the Extended Twisted Mass
  Collaboration: the hadronic vacuum polarisation, where I am the reference person of the Bern group in the
  blinded analysis \ref{pub:hvpupdate}, \ref{pub:g2etm}, and the hadronic light-by-light contribution, where I am the contact person for the pole contributions and the analysis tool
  chain \ref{pub:hlbl}.
  \item the calculation of leading isospin-breaking effects with the RM123 method, in particular the
  neutron-proton mass difference and the mass splittings of the $\Delta(1232)$ resonances, which was the
  subject of my doctoral thesis and one of the few theoretical determinations available in the literature
  \ref{pub:isospin}, \ref{pub:nptalk}.
  \item methods and software for lattice calculations: the extraction of multiple exponential signals from
  Euclidean correlators \ref{pub:extraction}, the auto-tuning of multigrid solver parameters in the hybrid Monte Carlo \ref{pub:multigrid}, and
  GPU implementations for gauge-ensemble generation \ref{pub:gpuensemble}.
\end{tightitemize}

\cvhead{Professional experience}

\begin{cvlist}
\item[01/04/2024 - today] \textbf{Postdoctoral researcher}, \textit{University of Bern}, Bern, Switzerland\\
Institute for Theoretical Physics and Albert Einstein Center for Fundamental Physics. Member of the Extended
Twisted Mass Collaboration (see Research activity above and Participation in international collaborations
below).\\
Supervisor: Prof. Urs Wenger.

\item[01/11/2021 - 31/03/2024] \textbf{Postdoctoral researcher}, \textit{University of Bonn}, Bonn, Germany\\
Helmholtz Institute for Radiation and Nuclear Physics (HISKP), Computational Physics group. Member of the
Extended Twisted Mass Collaboration (see Research activity above and Participation in international
collaborations below).\\
Supervisor: Prof. Carsten Urbach.

\item[01/11/2018 - 31/10/2021] \textbf{Doctoral researcher}, \textit{Roma Tre University}, Rome, Italy\\
Department of Mathematics and Physics (see Research activity above).\\
Supervisors: Prof. Vittorio Lubicz, Dr. Silvano Simula.
\end{cvlist}

\cvhead{Education}

\begin{cvlist}
\item[01/11/2018 - 31/10/2021] \textbf{PhD in Theoretical Physics}, \textit{Roma Tre University}, Rome, Italy\\
Thesis: \textit{The neutron--proton mass difference in Lattice QCD+QED}.\\
Supervisors: Vittorio Lubicz, Silvano Simula.\\
Date of submission: 31/10/2021. Date of defence: 22/04/2022.\\
Admitted first in the ranking of the national competitive admission examination.

\item[01/10/2016 - 23/10/2018] \textbf{Master's Degree in Theoretical Physics of Elementary Particles},
\textit{Roma Tre University}, Rome, Italy\\
Thesis: \textit{Optimization techniques in the lattice calculation of the hadronic contribution to the muon
anomaly}.\\
Supervisor: Prof. Cecilia Tarantino.\\
Final mark: 110/110 cum laude. GPA: 29.85/30. Date of defence: 23/10/2018.

\item[01/10/2013 - 21/07/2016] \textbf{Bachelor's Degree in Physics}, \textit{Roma Tre University}, Rome, Italy\\
Thesis: \textit{Astro-Archaeology of the Galactic Centre}.\\
Supervisor: Prof. Giorgio Matt.\\
Final mark: 110/110 cum laude. GPA: 28.84/30. Date of defence: 21/07/2016.
\end{cvlist}

\cvhead{Publications}

{\raggedright\noindent\textbf{ORCID:}~\href{https://orcid.org/0000-0002-6509-447X}{orcid.org/0000-0002-6509-447X}\quad
\textbf{Inspire:}~\href{https://inspirehep.net/authors/1755571}{inspirehep.net/authors/1755571}\quad
\textbf{GitHub:}~\href{https://github.com/simone-romiti}{github.com/simone-romiti}\par}

\cvsub{Peer-reviewed journals}

\begin{enumerate}[leftmargin=1.6em,labelsep=0.5em,label={[J\arabic*]},itemsep=2pt,topsep=3pt]
  \item \label{pub:pinn} \textbf{S. Romiti}, \textit{$\mathrm{SU}(N)$ lattice gauge theories with physics-informed neural
  networks}, Phys.\ Rev.\ D \textbf{113} (2026) 054511, arXiv:2510.26904 [hep-lat].

  \item T. Jakobs, M. Garofalo, T. Hartung, K. Jansen, P. Ludwig, J. Ostmeyer, \textbf{S. Romiti}, C. Urbach,
  \textit{A comprehensive stress test of truncated Hilbert space bases against Green's function Monte Carlo in
  $\mathrm{U}(1)$}, Eur.\ Phys.\ J.\ C \textbf{86} (2026) 637, arXiv:2510.27611 [hep-lat].

  \item T. Jakobs, M. Garofalo, T. Hartung, K. Jansen, J. Ostmeyer, \textbf{S. Romiti}, C. Urbach,
  \textit{Dynamics in Hamiltonian lattice gauge theory: approaching the continuum limit with partitionings of
  $\mathrm{SU}(2)$}, Eur.\ Phys.\ J.\ C \textbf{85} (2025) 1418, arXiv:2503.03397 [hep-lat].

  \item \label{pub:matching} C. F. Gro\ss{}, \textbf{S. Romiti}, L. Funcke, K. Jansen, A. Kan, S. K\"uhn, C. Urbach,
  \textit{Matching Lagrangian and Hamiltonian simulations in $(2+1)$-dimensional $\mathrm{U}(1)$ gauge theory},
  Eur.\ Phys.\ J.\ C \textbf{85} (2025) 1253, arXiv:2503.11480 [hep-lat].

  \item \label{pub:runningcoupling} A. Crippa, \textbf{S. Romiti}, L. Funcke, K. Jansen, S. K\"uhn, P. Stornati, C. Urbach,
  \textit{Towards determining the $(2+1)$-dimensional Quantum Electrodynamics running coupling with Monte Carlo
  and quantum computing methods}, Commun.\ Phys.\ \textbf{8} (2025) 367, arXiv:2404.17545 [hep-lat].

  \item R. Aliberti \textit{et al.} (Muon $g-2$ Theory Initiative), \textit{The anomalous magnetic moment of the
  muon in the Standard Model: an update}, Phys.\ Rept.\ \textbf{1143} (2025) 1--158, arXiv:2505.21476 [hep-ph].

  \item \label{pub:g2etm} Extended Twisted Mass Collaboration (C. Alexandrou \textit{et al.}), \textit{Strange and charm quark
  contributions to the muon anomalous magnetic moment in lattice QCD with twisted-mass fermions}, Phys.\ Rev.\ D
  \textbf{111} (2025) 054502, arXiv:2411.08852 [hep-lat].

  \item Extended Twisted Mass Collaboration (C. Alexandrou \textit{et al.}), \textit{Inclusive hadronic decay
  rate of the $\tau$ lepton from lattice QCD: the $\bar{u}s$ flavour channel and the Cabibbo angle}, Phys.\ Rev.\
  Lett.\ \textbf{132} (2024) 261901, arXiv:2403.05404 [hep-lat].

  \item \label{pub:digitization} \textbf{S. Romiti} and C. Urbach, \textit{Digitizing lattice gauge theories in the magnetic basis:
  reducing the breaking of the fundamental commutation relations}, Eur.\ Phys.\ J.\ C \textbf{84} (2024) 708,
  arXiv:2311.11928 [hep-lat].

  \item \label{pub:canonicalmomenta} T. Jakobs, M. Garofalo, T. Hartung, K. Jansen, J. Ostmeyer, D. Rolfes, \textbf{S. Romiti}, C. Urbach,
  \textit{Canonical momenta in digitized $\mathrm{SU}(2)$ lattice gauge theory: definition and free theory},
  Eur.\ Phys.\ J.\ C \textbf{83} (2023) 669, arXiv:2304.02322 [hep-lat].

  \item \label{pub:isospin} \textbf{S. Romiti}, \textit{Leading isospin breaking effects in nucleon and $\Delta$ masses},
  Rev.\ Mex.\ Fis.\ Suppl.\ \textbf{3} (2022) 0308030, arXiv:2203.04636 [hep-lat].

  \item \label{pub:extraction} \textbf{S. Romiti} and S. Simula, \textit{Extraction of multiple exponential signals from lattice
  correlation functions}, Phys.\ Rev.\ D \textbf{100} (2019) 054515, arXiv:1907.09926 [hep-lat].
\end{enumerate}

\cvsub{Conference proceedings}

\begin{enumerate}[leftmargin=1.6em,labelsep=0.5em,label={[C\arabic*]},itemsep=2pt,topsep=3pt]
  \item \label{pub:hvpupdate} ETM Collaboration (S. Bacchio \textit{et al.}, incl.\ \textbf{S. Romiti}), \textit{An update on the HVP contribution to $g_\mu-2$ in isoQCD from ETMC}, 42nd Int.\ Symp.\ on Lattice Field Theory (2026), arXiv:2603.03120 [hep-lat].
  \item \label{pub:hlbl} N. Kalntis, G. Kanwar, M. Petschlies, \textbf{S. Romiti}, U. Wenger, \textit{HLbL contribution to the muon $g-2$ using twisted-mass fermions at the physical point}, 42nd Int.\ Symp.\ on Lattice Field Theory (2025), arXiv:2512.21155 [hep-lat].
  \item \label{pub:multigrid} M. Garofalo, B. Kostrzewa, \textbf{S. Romiti}, A. Sen, \textit{Autotuning multigrid parameters in the HMC on different architectures}, PoS \textbf{LATTICE2024} (2025) 276.
  \item N. Kalntis, G. Kanwar, M. Petschlies, \textbf{S. Romiti}, U. Wenger, \textit{Hadronic light-by-light contribution to the muon $g-2$ using twisted-mass fermions}, PoS \textbf{LATTICE2024} (2025) 228, arXiv:2412.12320 [hep-lat].
  \item A. Evangelista \textit{et al.} (Extended Twisted Mass, incl.\ \textbf{S. Romiti}), \textit{Valence leading isospin breaking contributions to $a_\mu^{\mathrm{HVP\text{-}LO}}$}, PoS \textbf{LATTICE2024} (2025) 260, arXiv:2501.19350 [hep-lat].
  \item M. Garofalo, T. Hartung, T. Jakobs, K. Jansen, J. Ostmeyer, D. Rolfes, \textbf{S. Romiti}, C. Urbach, \textit{Testing the $\mathrm{SU}(2)$ lattice Hamiltonian built from $S_3$ partitionings}, PoS \textbf{LATTICE2023} (2024) 231, arXiv:2311.15926 [hep-lat].
  \item M. Garofalo, T. Hartung, K. Jansen, J. Ostmeyer, \textbf{S. Romiti}, C. Urbach, \textit{Defining canonical momenta for discretised $\mathrm{SU}(2)$ gauge fields}, PoS \textbf{LATTICE2022} (2023) 040, arXiv:2210.15547 [hep-lat].
  \item L. Funcke, C. F. Gro\ss{}, K. Jansen, S. K\"uhn, \textbf{S. Romiti}, C. Urbach, \textit{Hamiltonian limit of lattice QED in $2+1$ dimensions}, PoS \textbf{LATTICE2022} (2023) 292, arXiv:2212.09627 [hep-lat].
  \item \label{pub:gpuensemble} B. Kostrzewa, S. Bacchio, J. Finkenrath, M. Garofalo, F. Pittler, \textbf{S. Romiti}, C. Urbach (ETM), \textit{Twisted mass ensemble generation on GPU machines}, PoS \textbf{LATTICE2022} (2023) 340, arXiv:2212.06635 [hep-lat].
  \item \label{pub:nptalk} \textbf{S. Romiti}, \textit{The neutron--proton mass difference}, PoS \textbf{LATTICE2021} (2022) 543, arXiv:2202.01613 [hep-lat].
\end{enumerate}

\cvsub{Published software}

\begin{enumerate}[leftmargin=1.6em,labelsep=0.5em,label={[\arabic*]},itemsep=3pt,topsep=3pt]
  \item \textbf{S. Romiti}, \textit{\href{https://github.com/simone-romiti/adiabatic_PINNs-suN}{\texttt{adiabatic\_PINNs-suN}}}: creator. Standalone implementation accompanying the adiabatic-PINN paper. It reproduces the published results.
  \item \textbf{S. Romiti} and C. Urbach, \textit{\texttt{su2\_DJT}} (2023), Zenodo, \href{https://doi.org/10.5281/zenodo.10143897}{doi:10.5281/zenodo.10143897}: creator. Archived and citable implementation accompanying the $\mathrm{SU}(2)$ digitisation paper.
  \item \textbf{S. Romiti}, \textit{\href{https://github.com/simone-romiti/lattice_data_tools}{\texttt{lattice\_data\_tools}}}: creator and maintainer. Actively maintained toolkit for the production and analysis of lattice data, L-CNN, canonical momenta, solution of Kogut-Susskind Hamiltonian with PINNs, etc.
  \item \textit{\href{https://github.com/HISKP-LQCD/su2}{\texttt{su2}}}: main maintainer. Monte Carlo library for $\mathrm{U}(1)$, $\mathrm{SU}(2)$ and $\mathrm{SU}(3)$ in $d=2,3,4$ dimensions, used for several of the publications above and beyond my own group.
  \item \textit{\href{https://github.com/etmc/tmLQCD}{\texttt{tmLQCD}}}: contributor. Community library for twisted-mass lattice QCD, used to generate the ETM collaboration ensembles.
\end{enumerate}

\cvhead{Talks at international conferences, workshops and seminars}

\begin{cvlist}
\item[30/09-02/10/2026] \textbf{Invited speaker.} \href{https://agenda.infn.it/event/53069/}{StaND4LQ 2026}, \textit{State
of the Art \& New Directions for Lattice Gauge Theory and Quantum Simulation}, University of Pisa, Italy.
First edition of this workshop, organised entirely by PhD students, bringing together researchers in
lattice gauge theories, quantum simulation and quantum many-body physics.

\item[03-07/08/2026] \textbf{Invited speaker.} \href{https://indico.phys.uconn.edu/event/3/}{9th Plenary
Workshop of the Muon $g-2$ Theory Initiative}, University of Connecticut, Storrs, United States.
\href{https://indico.phys.uconn.edu/event/3/contributions/104/}{Talk} on the update of
$a_\mu^{\mathrm{HVP,\,isoQCD}}$ from the ETM collaboration.

\item[26/07-01/08/2026] \href{https://indico.global/event/16565/}{Lattice 2026}, the 43rd International
Symposium on Lattice Field Theory, University of Maryland, College Park, United States.
\href{https://indico.global/event/16565/contributions/156143/}{Talk} on
$\mathrm{SU}(N)$ lattice gauge theories with physics-informed neural networks.

\item[08-13/03/2026] \textbf{Invited speaker.} \href{https://indico.cern.ch/category/376/}{CERN Lattice
Seminar}, CERN, Geneva, Switzerland. \href{https://indico.cern.ch/event/1641317/}{Talk} on the adiabatic PINN algorithm
for $\mathrm{SU}(N)$ lattice gauge theories.

\item[25-27/03/2026] \href{https://indico.global/event/16475/}{ETM collaboration meeting}, Cyprus Institute,
Nicosia, Cyprus. Progress of the Bern group on the HLbL contribution to $(g-2)_\mu$.

\item[01-05/12/2025] \href{https://indico.ectstar.eu/event/248/}{ECT* workshop} \textit{Multi-canonical methods
and lattice field theory}, Trento, Italy.
\href{https://indico.ectstar.eu/event/248/contributions/6627/}{Talk} on the adiabatic learning of the
coupling dependence of eigenfunctions and energies.

\item[08-12/09/2025] \href{https://indico.ijclab.in2p3.fr/event/11652/}{8th Plenary Workshop of the Muon $g-2$
Theory Initiative}, IJCLab, Orsay, France.
\href{https://indico.ijclab.in2p3.fr/event/11652/contributions/39079/}{Talk} on behalf of the ETM
collaboration on the HLbL contribution.

\item[01-05/09/2025] \href{https://indico.ectstar.eu/event/242/}{ECT* workshop} \textit{Hamiltonian Lattice
Gauge Theories: Status, Novel Developments and Applications}, Trento, Italy. Closing
\href{https://indico.ectstar.eu/event/242/contributions/5858/}{summary talk} of the workshop from the
organisers, with C. Urbach.

\item[03-07/03/2025] \textbf{Invited speaker.} \href{https://indico.ectstar.eu/event/229/}{ECT*
workshop} \textit{Scale Setting: Precision Lattice QCD for Particle and Nuclear Physics}, Trento, Italy.
\href{https://indico.ectstar.eu/event/229/contributions/5308/}{Talk} on the sub-permille determination of the
lattice spacing in isoQCD within the ETM collaboration.

\item[28/07-03/08/2024] \href{https://conference.ippp.dur.ac.uk/event/1265/}{Lattice 2024}, the 41st
International Symposium on Lattice Field Theory, Liverpool, United Kingdom.
\href{https://conference.ippp.dur.ac.uk/event/1265/contributions/7618/}{Talk} on the Nested Sampling algorithm for
$\mathrm{SU}(N)$ lattice gauge theories (the first of three coordinated talks in that session).

\item[31/07-04/08/2023] \href{https://indico.fnal.gov/event/57249/}{Lattice 2023}, the 40th International
Symposium on Lattice Field Theory, Fermilab, Batavia, United States.
\href{https://indico.fnal.gov/event/57249/contributions/268337/}{Talk} on the digitisation of lattice gauge
theories in the magnetic basis.

\item[20-24/03/2023] \href{https://indico.ectstar.eu/event/164/}{ECT* Gradient Flow Workshop}, Trento, Italy.
\href{https://indico.ectstar.eu/event/164/contributions/3536/}{Talk} on the gradient-flow determination of the
renormalised anisotropy on anisotropic lattices.

\item[08-13/08/2022] \href{https://indico.hiskp.uni-bonn.de/event/40/}{Lattice 2022}, the 39th International
Symposium on Lattice Field Theory, Bonn, Germany.
\href{https://indico.hiskp.uni-bonn.de/event/40/contributions/704/}{Talk} on the Hamiltonian limit of lattice
gauge theories.

\item[26-30/07/2021] \href{https://indico.cern.ch/event/1006302/}{Lattice 2021}, the 38th International
Symposium on Lattice Field Theory, MIT, United States (online).
\href{https://indico.cern.ch/event/1006302/contributions/4373299/}{Talk} on the lattice determination of the
neutron--proton mass difference.

\item[26-30/07/2021] \href{https://indico.nucleares.unam.mx/event/1541/}{Hadron 2021}, the 19th International Conference on Hadron Spectroscopy and Structure
(online). \href{https://indico.nucleares.unam.mx/event/1541/session/28/contribution/226}{Talk} on the mass
splittings of the nucleons and of the $\Delta(1232)$ resonances.
\end{cvlist}

\cvhead{Attendance of international workshops and schools}

\begin{cvlist}
\item[06-17/07/2026] \href{https://indico.cern.ch/event/1616325/}{Lattice@CERN 2026}, \textit{Improved estimators for lattice QCD correlation functions},
CERN Theory Department, Geneva, Switzerland.
\item[12-14/02/2025] \href{https://indico.psi.ch/event/16749/}{Z\"urich Gradient Flow Workshop 2025},
University of Zurich, Switzerland.
\end{cvlist}

\cvhead{Awards and grants}

\begin{cvlist}
\item[2025 - 2026] \textbf{Person of Contact} for Computing-time allocation of $22.90\times10^{6}$ core-hours on the JUPITER Booster
module, awarded by the Gauss Centre for Supercomputing (GCS).

\item[2025 - 2026] \textbf{Principal investigator} of a computing-time allocation of $0.20\times10^{6}$
core-hours on ALPS at CSCS, the Swiss National Supercomputing Centre, corresponding to
$0.80\times10^{6}$ GPU node-hours, for the continuum-limit calculation of the $\eta,\eta'$ transition form
factors on the ETM ensembles.

\item[2026] Winner of a CHF 5{,}000 grant for the organisation of the \textit{Swiss Lattice Meeting}, the first
national gathering for lattice gauge theory in Switzerland, to be held on 14-15/01/2027.

\item[2026] Winner of a CHF 3{,}000 grant from the \textit{Open Round 2026} programme of the University of Bern,
for a visiting research period at TU Wien, Vienna, Austria, from 15/05/2026 to 15/06/2026.

\item[2013 - 2016] Winner of a EUR 3{,}000 merit scholarship awarded by Roma Tre University for academic
excellence during the Bachelor's degree.
\end{cvlist}

\cvhead{Event organisation}

\begin{cvlist}
\item[14-15/01/2027] \textbf{Organiser} of the \textit{Swiss Lattice Meeting}, one of four organisers, all
postdoctoral researchers. First edition of the national meeting for lattice gauge theory in Switzerland,
established as a recurring event. I secured the CHF 5{,}000 of funding for it.

\item[01-05/09/2025] \textbf{Leading organiser} of the \href{https://indico.ectstar.eu/event/242/}{ECT*
workshop} \textit{Hamiltonian Lattice Gauge Theories: Status, Novel Developments and Applications}, Trento, Italy.
I proposed the workshop and led its organisation over the week.

\item[20-24/01/2025] \textbf{Organising support} for the
\href{https://indico.global/event/9559/}{\textit{Workshop on the sign problem in QCD and beyond}} (SIGN25),
Institute for Theoretical Physics, University of Bern, Switzerland.

\item[2022 - 2024] \textbf{Co-organiser}, with M. Mai, of the
\href{https://indico.hiskp.uni-bonn.de/category/35/}{weekly physics seminars} of HISKP, University of Bonn, Germany:
invitation of the speakers, chairing of the sessions and broadcasting.

\item[08-13/08/2022] \textbf{Organising support} for
\href{https://indico.hiskp.uni-bonn.de/event/40/}{Lattice 2022}, the 39th International Symposium on Lattice
Field Theory, Bonn, Germany, as one of the local organisers, including the audio and video broadcasting of the
scientific programme.
\end{cvlist}

\cvhead{Participation in international collaborations}

\begin{cvlist}
\item[01/04/2024 - today] Member of the Extended Twisted Mass (ETM) Collaboration, as reference person of the
Bern group for the blinded analysis of the hadronic vacuum polarisation and contact person for the hadronic
light-by-light contribution.
\item[01/11/2021 - 31/03/2024] Member of the Extended Twisted Mass (ETM) Collaboration within the Bonn group.
\item[2025 - today] Contributor to the Muon $g-2$ Theory Initiative.
\end{cvlist}

\cvhead{Supervision and teaching}

\noindent\hspace{0.5cm}\begin{minipage}{\dimexpr\textwidth-0.5cm\relax}
\textit{Supervision.} Within my postdoctoral positions I have guided graduate students on both the theoretical
understanding and the technical execution of their projects, through regular meetings and progress reviews
against their milestones.
\end{minipage}

\begin{cvlist}
\item[2024 - today] University of Bern, Switzerland: three PhD students (N. Kalntis, J. F. de la Garza,
L. N. Filipovic).
\item[2022 - 2024] University of Bonn, Germany: two MSc and PhD students (C. F. Groß, T. Jakobs, O. Luniachek).
\end{cvlist}

\noindent\hspace{0.5cm}\begin{minipage}{\dimexpr\textwidth-0.5cm\relax}
\textit{Teaching.} Tutorials and assistantships over eight semesters, for six university courses. The teaching material of several of
them is maintained in open repositories on GitHub.
\end{minipage}

\begin{cvlist}
\item[01/10/2023 - 01/02/2024] \textit{Computational Physics} (Prof. C. Urbach), University of Bonn, Germany.
\item[01/04/2022 - 31/07/2022] \textit{Quantum Computing} (Prof. L. Funcke), University of Bonn, Germany.
\item[01/10/2021 - 28/02/2022] \textit{Quantum Computing} (Prof. C. Urbach), University of Bonn, Germany.
\item[01/10/2020 - 28/02/2021] \textit{Laboratory of Programming and Calculus} (Prof. S. Bussino), Roma Tre
University, Italy.
\item[01/10/2019 - 28/02/2020] \textit{Mathematical Analysis 1} (Prof. U. Bessi), Roma Tre University, Italy.
\item[01/10/2020 - 28/02/2021] \textit{Mathematical Analysis 1} (Prof. L. Chierchia), Roma Tre University, Italy.
\item[01/10/2020 - 28/02/2021] \textit{Fundamentals of Mathematics 2} (Prof. L. Tedeschini Lalli), Department
of Architecture, Roma Tre University, Italy, as assistant of the lecturer.
\item[01/02/2020 - 30/06/2020] \textit{Physics 1} (Prof. W. Plastino), Roma Tre University, Italy.
\end{cvlist}

\cvhead{Outreach}

\begin{cvlist}
\item[15/02/2019] \textbf{Co-organiser and presenter} of the public event
\href{https://matematicafisica.uniroma3.it/terza-missione/iniziative-per-tutti/occhi-su-marte/}{\textit{Occhi su
Marte}} (``Eyes on Mars''), Department of Mathematics and Physics, Roma Tre University, Rome, Italy. The event
consisted of seminars, hands-on activities for children and telescope observations for the general public. I
contributed to the organisation of the whole event and presented a live demonstration of the Barkhausen effect.
\end{cvlist}

\cvhead{Peer review}

\begin{cvlist}
\item[2023, 2024] Referee for \textit{Physical Review D}, one manuscript in each year.
\end{cvlist}

\cvhead{Technical skills}

\begin{cvlist}[2.6cm]
\item[Languages] C, C++, Python, Bash, R, SQL, \LaTeX{}, Quarto.
\item[HPC] CUDA, MPI, OpenMP, SLURM, GNU/Linux clusters.
\item[Scientific ML] PyTorch, NumPy, SciPy, SymPy.
\item[Development] Git, GitHub Actions, continuous integration, Docker.
\item[Methods] Monte Carlo and hybrid Monte Carlo, Bayesian inference and model averaging, physics-informed
neural networks, variational autoencoders, diffusion models.
\end{cvlist}

\cvhead{Spoken languages}

\begin{cvlist}[2.6cm]
\item[Italian] Mother tongue
\item[English] Fluent
\item[German] Intermediate
\end{cvlist}